% Source: https://artofproblemsolving.com/community/c3h2360247_math_advent_calendar
% Stable identity: aops-topic-2360247-q10
Consider three cases.
\begin{itemize}
\item $x\ge6$: the inequality is $5x(x+3)+2\ge x-6$, i.e. $5x^2+14x+8\ge0$, which is true for $x\ge6$.
\item $-3\le x<6$: it is $5x(x+3)+2\ge6-x$, i.e. $5x^2+16x-4\ge0$. The roots of $5x^2+16x-4$ are $\dfrac{-8\pm2\sqrt{21}}{5}$, so this holds for $x\le\dfrac{-8-2\sqrt{21}}{5}$ or $x\ge\dfrac{-8+2\sqrt{21}}{5}$. Since $\dfrac{-8-2\sqrt{21}}{5}\approx-3.43<-3$, in this case the solutions are $\dfrac{-8+2\sqrt{21}}{5}\le x<6$.
\item $x<-3$: here $|x+3|=-x-3$ and the inequality is $-5x(x+3)+2\ge6-x$, i.e. $5x^2+14x+4\le0$, which holds only for $\dfrac{-7-\sqrt{29}}{5}\le x\le\dfrac{-7+\sqrt{29}}{5}$, approximately $-2.48\le x\le-0.32$. None of these $x$ is less than $-3$.
\end{itemize}
So the solution set is $x\ge\dfrac{2\sqrt{21}-8}{5}$, and the least solution is
\[\boxed{x=\dfrac{2\sqrt{21}-8}{5}}\approx0.233.\]
